\chapter{Introduction}
\label{chap:introduction}

\section{Background and Motivation}

Probabilistic Risk Assessment (PRA) is a structured, model based approach used to understand risk in complex engineered systems such as nuclear power plants. In practice, PRA (1) identifies possible initiating events and failure scenarios, (2) estimates how likely those scenarios are to occur (e.g., frequencies or probabilities), (3) evaluates their potential consequences, and (4) explicitly represents and propagates the uncertainties in the underlying models and data. By linking plant systems, operator actions, and accident progressions through models such as event trees and fault trees, probabilistic risk assessment (PRA) generates quantitative risk metrics that characterize potential failure scenarios, the mechanisms by which they may occur, their estimated frequencies and consequences, and the confidence associated with those estimates.

PRA insights are usually combined with deterministic analyses, engineering judgment, and explicit safety margins to support risk-informed decision-making. This integration helps utilities and regulators prioritize attention and resources toward the most safety significant systems, components, and operational actions, while maintaining the conservative principles required in nuclear safety.

A central part of PRA is quantification: turning large logical models and reliability data into numerical results. Many commonly used quantification activities are computationally intensive, especially when performed repeatedly for configuration changes or model updates. Examples of such demanding tasks include computing (i) system failure probabilities from fault trees, (ii) event sequence frequencies from event trees, (iii) safety integrity level related metrics where applicable, (iv) common cause failure impacts, (v) importance measures used to identify risk significant components and operator actions, and (vi) the propagation of parameter uncertainty and (vii) sensitivity studies. These are not the only quantifications performed in PRA, but they represent a set of core calculations that frequently dominate runtime and are central to day-to-day risk studies.

The computational burden increases further as the scope of PRA expands to multi-hazard and multi-unit scenarios. External hazards, coupled dependencies, shared systems, and correlated failures can significantly increase the model size, producing very complex logic models.

For the nuclear energy sector, the ability to quantify these complex PRA models is not just a technical convenience, it is also linked to safety assurance and regulatory compliance. Plants must be able to perform thorough analyses, update risk models, support audits and reviews, and, deliver results on timelines that align with operational needs.

Despite the importance of these analyses, the current PRA software ecosystem remains difficult to access, difficult to integrate across platforms, and difficult to scale to the computational demands of complex PRA models. PRA tools are often inaccessible due to licensing issues and costs, and they are incompatible across workflows because the ecosystem is fragmented. This fragmentation is driven by several practical realities: nuclear plant data and models are often proprietary; tool developers have adopted different architectural and methodological philosophies (for example, differences in event tree linking versus fault tree linking approaches); and there is no standard data format that consistently represents PRA inputs, intermediate data structures, and results across tools. As a result, each platform tends to define its own data structures (often layered on top of general purpose formats such as XML, JSON etc.), and transferring a model from one tool to another typically requires tool specific knowledge. Conversions can be time consuming, error prone, and sometimes inherently lossy when data fields do not map properly. Operating system dependencies and vendor lock-in incentives further reinforce these issues.

Performance and scalability limitations compound the problem. Many existing tools are designed primarily for single user desktop workflows, offering serial execution or limited parallelism, often exposed through outdated GUIs or script oriented interfaces. For complex models, full quantification can take hours or even days, and practitioners may have limited visibility into whether a long running job is progressing normally or has stalled. When runs fail or appear stuck after extensive runtime, analysts may be forced to restart calculations, wasting time and compute resources. These constraints make it difficult to support time sensitive use cases (for example, near real-time evaluations during plant operation) and frequently encourage the use of truncation limits to keep calculations tractable. In practice, truncation thresholds are often selected based on experience and engineering conventions, but there is limited universal guidance, and the tradeoff between runtime and completeness can be difficult to manage transparently.

However, recent advances in open source containerization, message queueing, and microservices architectures create a clear opportunity to modernize PRA quantification workflows. A domain specific, distributed queueing and execution service can provide a unified API gateway that is operating system agnostic, supports automated deployments, and coordinates execution across modern heterogeneous hardware (GPUs, multi-core CPUs, FPGAs etc). Such a platform can also provide observability, traceability, and managerial control features that are important in regulated environments. The same design should support both cluster based deployments and a single node distributed memory based deployment, where each CPU core and available GPU can operate as an independent worker process with isolated resources (memory, compute allocation, and I/O), enabling distributed execution even without a multi machine cluster.

This work pursues that direction by designing and evaluating a distributed networked quantification platform for PRA. The study focuses on how such a system affects key operational and computational outcomes such as scalability, throughput, latency, and interoperability aiming to reduce analysis time, improve reliability and observability of long running quantifications, and lower integration barriers for PRA practitioners, tool developers, and organizations that must collaborate on complex risk models.

\section{Objectives}
This research work aims to accomplish the following objectives:
\begin{enumerate}
    \item Design a domain specific, operating system agnostic microservice that exposes a unified REST API gateway for PRA quantification and supports multi user, asynchronous submission.
    \item Define a common job specification and interoperability layer that decouple model submission from solver specific formats and interfaces.
    \item Integrate multiple PRA solvers by creating tool specific endpoints, demonstrating how both Windows and Linux based solvers can be supported.
    %\item Design and implement a priority-aware scheduling mechanism for a distributed queueing system that reduces latency for high-priority tasks (e.g., uncertainty quantification) while maintaining high throughput for mixed workloads under heavy load.
    \item Develop worker runtimes for CPUs and GPUs that scale across on-premise clusters or across distributed memory based deployments by treating each CPU core and each GPU on a node as an independent worker.
    \item Provide preprocessing, processing, and postprocessing workflows to enable distributed quantification of large models, including very large event trees.
    \item Leverage open source container orchestration to automate most deployment, scaling, and life cycle management tasks while preserving manual control where needed.
    \item Instrument the platform with standardized telemetry for latency, throughput, and error rates to support monitoring.
    \item Conduct a comprehensive evaluation that compares serialized, parallel, and distributed execution, quantifies speedup across models and calculation types, and characterizes scaling with the number of workers.
    \item Enforce role managed access and isolated execution to accommodate proprietary plant data and license restrictions.
    %assesses priority handling under mixed workloads,
    %and analyzes overheads for small, medium, and large models.
    \item Provide comprehensive documentation including implementation details, deployment configurations, experiment scripts, and functional and benchmark tests to enable reproducibility, guide users in deploying the system, and support integration of new solvers into the distributed platform.
\end{enumerate}

\section{Thesis Organization}
This thesis is structured as follows. Chapter 2 reviews the current PRA software landscape, identifies limitations in current workflows, and surveys related work in distributed computing for PRA. Building on this foundation, chapter 3 establishes the motivation for a distributed approach to PRA quantification and introduces the relevant distributed systems theory, architectural patterns, and design requirements. Chapter 4 then presents the overall system design, detailing its architecture, components, data flow, and key considerations for scalability and interoperability. The implementation of this design is described in chapter 5, covering the technology stack, task execution mechanisms, and deployment strategies. Chapter 6 presents the experimental methodology and results from a case study on MHTGR models and summarizes the key findings. These results are then discussed, including a comparison with a state-of-the-art tool supported by the NRC and an analysis of the distributed system's strengths and limitations. Finally, Chapter 7 concludes by summarizing the contributions and outlining directions for future work. Supplementary material, including source code, experiment scripts, and user documentation is provided in the Appendix.

%%\section{Problem Statement}
%%PRA practitioners lack a publicly available, locally deployable platform that scales quantification across heterogeneous resources while preserving tool interoperability, data governance, and reproducibility. In particular:
%%\begin{enumerate}
%%    \item Tool heterogeneity: PRA solvers use different model formats, assumptions, and execution interfaces, complicating interoperability and comparative studies.
%%    \item Limited scalability: Many tools prioritize serial execution or ad hoc parallelization and do not provide a distributed, multi tenant service for submitting and scheduling large batches of jobs.
%%    \item Operational constraints: Operating system dependencies and desktop centric deployments hinder reproducible and portable execution across environments.
%%    \item Governance and licensing: Proprietary plant data and tool licenses necessitate on-premises, role managed execution with strict isolation and auditability.
%%    \item Workload management: Mixed workloads require priority aware scheduling to control latency and fair distribution of computational resources under load.
%%    \item Resource heterogeneity: Practical deployments must utilize CPUs and GPUs efficiently, including single node scenarios where each core and the integrated GPU serve as independent workers.
%%   \item Distributed model decomposition and aggregation: Users need the capability to partition large PRA models during preprocessing, execute the partitions concurrently across distributed workers, and deterministically aggregate the partial results during postprocessing to fully exploit the scalability of the distributed system.
%%\end{enumerate}
%%The research problem is to design, implement, and evaluate a distributed queueing system that addresses these constraints for PRA quantification without relying on external, closed infrastructures, and without forcing analysts to rework existing models or tools.